En \ocaml, une procédure peut appeler une autre procédure précédemment 
déclarée. Une procédure peut également s'appeler elle-même. On dit
alors qu'elle est {\em récursive}. Il faut alors la déclarer avec
{\mcl let rec} au lieu de {\mcl let} (de l'anglais {\em \underline{rec}ursive}). Les procédures récursives \index{Récursif}sont utiles
pour faire des dessins comme ceux de la figure \ref{fig:desrec} (page
\pageref{fig:desrec}). Ces dessins récursifs ont la particularité de se
retrouver indéfiniment en réduction à l'intérieur d'eux-mêmes (exactement
comme l'image infinie obtenue dans deux miroirs se faisant face). Comme en pratique il n'est 
pas possible de dessiner un trait plus petit qu'un pixel de l'écran, le tracé s'arrête 
pour les dessins très petits.
\begin{figure}[htbp]
 \includevspacePostScript[18.5cm]{DessinsRecursifs}{0.95}{0.95}
 \caption{Dessins récursifs}\label{fig:desrec}
\end{figure}

\section{Récursivité simple}

La structure des dessins récursifs simples se comprend par {\em récurrence} \index{Récurrence}sur la taille du dessin.
Il faut d'abord comprendre la {\em base\/}, c'est-à-dire comment faire un dessin de taille 0 (ou 1). 
Ensuite il faut comprendre l'{\em induction\/}, c'est-à-dire 
comment construire un dessin de taille $n$ en utilisant des dessins identiques de tailles $n-1$, 
$n-2$, \ldots strictement inférieures à $n$.  

Prenons l'exemple des pavés récursifs
de la figure \ref{fig:pavec} (page \pageref{fig:pavec}).
\begin{figure}[htbp]
 \includevspacePostScript[6.8cm]{PavesRecursifs}{0.95}{0.95}
 \caption{Pavés récursifs}\label{fig:pavec}
\end{figure}
Un pavé de taille 0 est réduit à un point (plus précisément un pixel de l'écran).
Un pavé de taille $n$ est constitué d'un pavé de taille $n - 1$
entouré d'un carré de côté $n$ fois la taille du quadrillage, comme expliqué à la
figure \ref{fig:indpavec} (page \pageref{fig:indpavec}).
\begin{figure}[htbp]
 \includevspacePostScript[8.3cm]{InductionPavesRecursifs}{0.95}{0.95}
 \caption{Règle de construction récurrente des pavés récursifs}\label{fig:indpavec}
\end{figure}
La procédure récursive {\idt pave} de dessin d'un pavé récursif reçoit en paramètre la taille {\idt n}
du pavé. Si {\idt n} est nul, elle dessine un point. Si {\idt n} est strictement positif alors
elle dessine un carré de côté {\idt n} puis s'appelle récursivement pour dessiner un
pavé de taille {\idt n - 1}.
\ \\
\input ocaml/PavesRecursifs.tex
\noindent Le mot-clé {\mcl rec} indique au compilateur que le nom {\idt pave} peut être
utilisé dans le corps de la procédure {\idt pave} elle-même. Si on l'oublie, le compilateur
signale `{\tt Unbound value pave}' (valeur {\idt pave} inconnue) en désignant l'appel
{\idt pave (n - 1)} et suggère `{\tt If this is a recursive definition, you should add the rec
keyword}' (s'il s'agit d'une définition récursive, il faut ajouter le mot-clé {\mcl rec}).


Le dessin des pavés récursifs peut se faire {\em itérativement}, c'est-à-dire en utilisant 
une boucle ``for'' ou ``while''. Ce serait plus difficile avec des dessins comme le tableau d'un tournoi
de tennis représenté à la figure \ref{fig:tabtour} ci-dessous.
\begin{figure}[htbp]
 \includevspacePostScript[4.3cm]{Tournoi1-2-4-8-16}{0.95}{0.95}%
 \caption{Règle de construction récurrente d'un tableau\label{fig:tabtour}}
 \centering{de tournoi de tennis}
\end{figure}

Un tableau de tournoi de tennis ne comprenant qu'un seul inscrit comporte un trait horizontal
prévu pour écrire le nom de cet inscrit qui est également le vainqueur. Un tableau de tournoi de 
tennis comprenant {\idt n} ins\-crits s'obtient en traçant un trait horizontal
prévu pour écrire le nom du vainqueur en face de la pointe d'une accolade reliant les 
sous-tableaux servant à désigner les finalistes. Chacun des sous-tableaux contient à peu près
le même nombre de joueurs c'est-à-dire respectivement {\idt n / 2} et
{\idt n - n / 2} joueurs (rappelons que {\idt /} est la division entière). La procédure récursive de dessin du tableau d'un 
tournoi de {\idt n} joueurs est donnée ci-dessous et illustrée dans le cas {\idt n = 11}.

 {
 %\vspace{\baselineskip}
 \vbox to 0pt{\vspace*{3cm}\includevspacePostScript[3cm]{Tournoi11}{0.95}{0.95}}%
 \input ocaml/TableauDeTournoi.tex
 }
\section{Invariants}
\index{Invariant}Quand on écrit une procédure récursive il est important de réfléchir à l'{\em in\-va\-riant\/} 
de cette procédure c'est-à-dire aux hypothèses que doivent satisfaire les paramètres
effectifs de la procédure et le \robot\ pour que le dessin soit correct. Prenons l'exemple d'un
arbre binaire (voir la figure \ref{fig:arbin} ci-dessous), une structure de base de la
programmation, que nous avons déjà rencontrée sous la forme du tableau de tournoi de tennis. 
\begin{figure}[htbp]
 \includevspacePostScript[6.6cm]{ArbreBinaire}{0.95}{0.95}
 \caption{Règle de construction récurrente d'un arbre binaire}\label{fig:arbin}
\end{figure}

\noindent La procédure de dessin d'un arbre binaire est la suivante~:

\input ocaml/ArbreBinaire.tex

Le premier invariant est que le paramètre {\idt l} de la procédure {\idt arbre} doit être 
supérieur ou égal à 1. C'est
vrai pour l'appel principal puisque {\idt l} $=$ 32. Un appel quelconque
ne fait des appels supplémentaires que si {\idt l} $>$ 1. Comme 
{\idt l} $\geq 2$, on a
$(${\idt l / 2}$) \geq 1$ et l'invariant est vrai quand on appelle récursivement la 
procédure pour dessiner le sous-arbre gauche. Comme 
{\idt l} $= (2 \ast (${\idt l / 2}$))$ 
si {\idt l} est pair et 
{\idt l} $= (2 \ast (${\idt l / 2}$)) + 1$
si {\idt l} est impair, on a
{\idt l} $\geq (2 \ast (${\idt l / 2}$))$ donc 
$(${\idt l} $- (${\idt l / 2}$)) \geq$ $(${\idt l / 2}$)$ de sorte que
$(${\idt l} $- (${\idt l / 2}$)) \geq 1$
puisque $(${\idt l / 2}$) \geq 1$. L'invariant est donc également vrai quand on appelle 
récursivement la 
procédure {\idt arbre} pour dessiner le sous-arbre droit. Comme l'invariant est vrai pour l'appel principal
et qu'il reste vrai pour tous les appels suivants, il est toujours vrai.

On vérifierait de la même façon l'invariant concernant le \robot, à savoir qu'avant et après
le dessin de l'arbre binaire, le \robot\ se trouve placé à la racine de l'arbre, orienté vers le nord avec son
crayon baissé.

\section{Non-terminaison}
Quand on utilise des procédures récursives, il peut arriver qu'elles s'appellent les unes les autres
sans fin, exactement comme dans le cas de boucles ``while'' dont l'exécution ne se termine pas. L'exemple le
plus simple est celui de la procédure {\idt boucler} ci-dessous.

\label{prog:NonTerminaison}
\input ocaml/NonTerminaison.tex
Un appel de {\idt boucler ()} provoque immédiatement un nouvel appel de {\idt bou\-cler ()}, et ainsi
de suite indéfiniment. Comme cet appel est la dernière chose que fait la procédure (on parle
d'{\em appel terminal}), le compilateur \ocaml\ le remplace par un simple retour au début de la
procédure~: le programme tourne alors sans fin, exactement comme une boucle ``while'' qui ne se
termine pas, et il faut l'arrêter en tapant \framebox{ctrl}~\framebox{C} dans le terminal. Si l'appel
récursif n'est pas terminal (par exemple si l'on écrit {\idt boucler ();\ av ()} dans le corps
de la procédure), l'ordinateur utilise un peu de sa mémoire à chaque appel pour se souvenir de ce
qui reste à faire ensuite. Il finit par se produire une erreur quand il n'y a plus
suffisamment de mémoire disponible pour faire l'appel suivant. Cette erreur est signalée 
par le message `{\tt Fatal error:\ exception Stack\_overflow}' (débordement de la pile) avant l'arrêt
de l'exécution du programme. En général 
l'erreur est moins flagrante et ne se produit que pour certaines valeurs des paramètres.

\section{Récursivité mutuelle}
Considérons les serpentines de la figure \ref{fig:serprec} (page \pageref{fig:serprec}). Il y en a
de deux sortes.  Les serpentines 
{\idt E$\scriptscriptstyle n$} ont 
une hauteur $n$ et se dessinent à partir de leur coin inférieur ouest, en traçant un trait 
horizontal d'ouest en est puis un trait vertical pour passer à l'est de l'étage supérieur. 
Les serpentines {\idt W$\scriptscriptstyle n$} ont une hauteur $n$ et se dessinent 
à partir de leur coin inférieur est, en traçant un trait horizontal d'est en ouest puis un trait
vertical pour passer à l'ouest de l'étage supérieur.  
\begin{figure}[htbp]
 \includevspacePostScript[4.4cm]{SerpentineRecursive}{0.95}{0.95}
 \caption{Règle de construction récurrente d'une serpentine}\label{fig:serprec}
\end{figure}

\noindent Les procédures {\idt serpentineE} et {\idt serpentineW} dans le programme ci-dessous
dessinent respectivement une serpentine {\idt E$\scriptscriptstyle n$} et 
{\idt W$\scriptscriptstyle n$}, leur paramètre formel {\idt n} ayant pour valeur effective
la hauteur $n$ de la serpentine à dessiner. 
\input ocaml/SerpentineRecursive.tex
\noindent Une serpentine
{\idt E$\scriptscriptstyle 0$} ou {\idt W$\scriptscriptstyle 0$} de hauteur nulle étant vide, 
les procédures {\idt serpentineE} et {\idt serpentineW} ne font rien quand la valeur $n$ de leur 
paramètre {\idt n} est négative ou nulle. 
Pour dessiner une serpentine {\idt E$\scriptscriptstyle n$} ou {\idt W$\scriptscriptstyle n$} dont la 
hauteur $n \geq 1$ est la valeur effective de leur paramètre formel {\idt n}, les procédures 
{\idt serpentineE} et {\idt serpentineW} dessinent la 
base puis au-dessus une serpentine {\idt W$\scriptscriptstyle n-1$} 
ou {\idt E$\scriptscriptstyle n-1$} dont  la hauteur est la valeur de {\idt n - 1}. Comme
la procédure {\idt serpentineE} appelle la procédure {\idt serpentineW} et que la procédure
{\idt serpentineW} appelle la procédure {\idt serpentineE}, on dit que ces procédures sont
{\em mutuellement récursives\/}.

\index{Identificateur}Quand on utilise des procédures mutuellement récursives en \ocaml, il faut
respecter la règle que {\em tout identificateur doit être déclaré avant d'être
utilisé}. Si l'on déclarait {\idt serpentineE} puis {\idt serpentineW} par deux
{\mcl let rec} séparés, le compilateur signalerait `{\tt Unbound value serpentineW}' dans le corps de
{\idt serpentineE}, puisque {\idt serpentineW} n'est pas encore déclarée. C'est pourquoi on
déclare les deux procédures {\em en même temps}~: la première avec {\mcl let rec} et la seconde
avec {\mcl and} (``et'' en anglais) à la place du {\mcl let}. Chacune peut alors appeler
l'autre (et elle-même).

De manière générale, les déclarations de procédures mutuellement
ré\-cur\-si\-ves ont  la forme suivante~:

\vfill

\index{Procédure}
\begin{center}
\begin{tabular}{@{}|p{12.5cm}|@{}} \hline
{\em Déclaration de procédures mutuellement récursives~:}\\ \hline \hline
\begin{minipage}[t]{12.5cm}
{\medskip

\mbox{\blc\ \ \ \cmt (*\ \small\rm Spécification de la première procédure\ \cmt *)}\par 
\mbox{\blc\ \ \ \mcl let rec\idt\ nomDeLaProcedure$_{1}$%
\footnote{\ Un nom de procédure est un identificateur.}%
\ paramFormel$_{1}^{1}$\footnote{\ Un paramètre formel est un identificateur.}\ paramFormel$_{1}^{2}$\ \dots\ =}\par 
\mbox{\blc\ \ \ \ \ \small\rm \ldots\ Instructions du corps de la procédure%
\footnote{\ Y compris les appels des procédures {\idt nom\-De\-La\-Pro\-ce\-du\-re$_{1}$} et
{\idt nom\-De\-La\-Pro\-ce\-du\-re$_{2}$}.}%
\footnote{\ Les expressions figurant dans ces instructions peu\-vent utiliser les paramètres formels
{\idt paramFormel$_{1}^{1}$}, {\idt paramFormel$_{1}^{2}$}, \ldots\ de la procédure
{\idt nom\-De\-La\-Pro\-ce\-du\-re$_{1}$}. }
 (séparées par des `$;$') \ldots}\par 

\medskip

\mbox{\blc\ \ \ \cmt (*\ \small\rm Spécification de la seconde procédure\ \cmt *)}\par 
\mbox{\blc\ \ \ \mcl and\idt\ nomDeLaProcedure$_{2}$%
\ paramFormel$_{2}^{1}$\ paramFormel$_{2}^{2}$\ \dots\ =}\par 
\mbox{\blc\ \ \ \ \ \small\rm \ldots\ Instructions du corps de la procédure%
\footnote{\ Les expressions figurant dans ces instructions peuvent utiliser les paramètres formels
{\idt paramFormel$_{2}^{1}$}, {\idt paramFormel$_{2}^{2}$}, \ldots\ de la procédure
{\idt nom\-De\-La\-Pro\-ce\-du\-re$_{2}$}.}
 (séparées par des `$;$') \ldots}\par 
\medskip}
\end{minipage}\\ \hline 
\end{tabular}
\end{center}

\noindent On peut de même déclarer ensemble trois procédures ou plus, en les séparant par des {\mcl and}.

\noindent \rule{137mm}{0.1mm} %---------------------------------------------
\begin{exercice}{L'exercice consiste à dessiner des {\em fractales} \cite{mandelbrot1}
\cite{mandelbrot2}, c'est-à-dire
des dessins que l'on retrouve à l'infini, en réduction à l'intérieur d'eux-mêmes.

\tiret  Écrire une procédure \ocaml\ récursive pour construire une courbe de Von Koch\/%
\footnote{\ Von Koch, mathématicien suédois, 1870--1924.}, comme
indiqué à la de la figure \ref{fig:segvonkoch} (page \pageref{fig:segvonkoch}). La courbe $K_{0}$
est un segment de droite.  La courbe $K_{1}$ est obtenue en divisant le segment   $K_{0}$ en trois segments 
égaux et en remplaçant celui du milieu par un triangle équilatéral sans base.
La courbe $K_{n}$ est obtenue de même en divisant en trois chacun des segments de la courbe
 $K_{n-1}$. Écrire un programme de dessin d'un flocon de Von Koch, obtenu en fragmentant ainsi vers l'extérieur
 les côtés d'un triangle équilatéral. Faire de même pour un anti-flocon de Von Koch pour lequel la 
 fragmentation se fait vers l'intérieur du triangle équilatéral initial. Utiliser diverses fragmentations
 d'un segment pour obtenir d'autres types de flocons comme ceux de la figure \ref{fig:flocvonkoch} 
 (page \pageref{fig:flocvonkoch}).

\tiret Programmer les dessins des motifs récursifs de la figure \ref{fig:desrec} 
(page \pageref{fig:desrec}).

\tiret Écrire un programme \ocaml\ de dessin des 
carrés (figure \ref{fig:carsierp}  page \pageref{fig:carsierp}),
triangles (figure \ref{fig:trisierp}  page \pageref{fig:trisierp}) et
pentagones (figure \ref{fig:pensierp}  page \pageref{fig:pensierp}) de Sierpinski\/%
\footnote{\ Waclaw Sierpinski, mathématicien polonais, 1882--1969.}.

\tiret Dessiner les végétaux de la figure \ref{fig:vege}  (page \pageref{fig:vege}).

\tiret Dessiner les carrés de Sedgewick \cite{sedgewick1} de la figure \ref{fig:sedgew}  (page \pageref{fig:sedgew}).

\tiret Dessiner la courbe du dragon de la figure \ref{fig:dragon}  (page \pageref{fig:dragon})
selon les règles de construction de
la figure \ref{fig:consdragon}  (page \pageref{fig:consdragon}). La courbe du dragon d'ordre $n$ est celle
que l'on observe en regardant par la tranche une feuille de papier qui a été repliée $n$ fois
sur elle-même puis dépliée à angles droits.

\tiret Dessiner la courbe de Hilbert\/%
\footnote{\ David Hilbert, mathématicien allemand, 1862--1943.}
de la figure \ref{fig:hilbert}  (page \pageref{fig:hilbert}) selon les règles de construction de
la figure \ref{fig:conshilbert}  (page \pageref{fig:conshilbert}).

\tiret Dessiner la courbe de Peano\/%
\footnote{\ Giuseppe Peano, mathématicien italien, 1858--1932.}
 de la figure \ref{fig:peano}  (page \pageref{fig:peano}) selon les règles de construction de
la figure \ref{fig:conspeano}  (page \pageref{fig:conspeano}).

\tiret Écrire une procédure \ocaml\ qui écrit en base 2 un nombre passé en paramètre
en remarquant que l'écriture binaire $\beta(n)$ d'un nombre $n$ est définie par~:
\begin{list}{$\scriptscriptstyle\bullet$}{
}
  \item $\beta(-n) = -\beta(n)$ si $-n < 0$$\:$;
  \item $\beta(0) = 0$$\:$;
 \item $\beta(1) = 1$$\:$;
 \item $\beta(n) = \beta(n $\ {\idt /}\ $ 2) \beta(n$\ {\mcl mod}\ $2)$ si $n > 1$.
\end{list}
Faire de même en bases 3, 8 et 16 (en utilisant les digits {\idt A}, {\idt B}, {\idt C}, {\idt D},
{\idt E}, {\idt F} pour dénoter respectivement les nombres 10, 11, 12, 13, 14, 15 et 
``{\idt ecrire} {\cha "d"}{\idt ; translater (1, 0)}'' 
pour écrire le digit \mbox{\cha d}).

\noindent $\Box$}\end{exercice}
\begin{figure}[htbp]
 \includevspacePostScript[5.1cm]{SegmentsVonKoch}{0.95}{0.95}
 \caption{Courbe de Von Koch}\label{fig:segvonkoch}
\end{figure}
\begin{figure}[htbp]
 \includevspacePostScript[2.4cm]{RegleDragon}{0.95}{0.95}
 \caption{Règles de construction de la courbe du dragon}\label{fig:consdragon}
\end{figure}
\begin{figure}[htbp]
 \includevspacePostScript[5.0cm]{RegleConsHilbert}{0.95}{0.95}
 \caption{Règles de construction des courbes de Hilbert}\label{fig:conshilbert}
\end{figure}
\begin{figure}[htbp]
 \includevspacePostScript[18.2cm]{FloconsVonKoch}{0.95}{0.95}
 \caption{Flocons de Von Koch}\label{fig:flocvonkoch}
\end{figure}
\begin{figure}[htbp]
 \includevspacePostScript[18.3cm]{CarreDeSierpinsky}{0.95}{0.95}
 \caption{Carrés de Sierpinski}\label{fig:carsierp}
\end{figure}
\begin{figure}[htbp]
 \includevspacePostScript[16.0cm]{TriangleDeSierpinsky}{0.95}{0.95}
 \caption{Triangles de Sierpinski}\label{fig:trisierp}
\end{figure}
\begin{figure}[htbp]
 \includevspacePostScript[16.0cm]{PentagoneSierpinski}{0.95}{0.95}
 \caption{Pentagones de Sierpinski}\label{fig:pensierp}
\end{figure}
\begin{figure}[htbp]
 \includevspacePostScript[12.5cm]{Vegetaux}{0.95}{0.95}
 \caption{Végétaux fractals}\label{fig:vege}
\end{figure}
\begin{figure}[htbp]
 \includevspacePostScript[14.0cm]{Sedgewick}{0.95}{0.95}
 \caption{Étoile fractale de Sedgewick}\label{fig:sedgew}
\end{figure}
\begin{figure}[htbp]
 \includevspacePostScript[7.2cm]{Dragon}{0.95}{0.95}
 \caption{Courbes du dragon d'ordre 1, 2, \ldots, 8, 10 et 12}\label{fig:dragon}
\end{figure}
\begin{figure}[htbp]
 \includevspacePostScript[17.5cm]{Hilbert}{0.95}{0.95}
 \caption{Courbes de Hilbert d'ordres 1 à 6}\label{fig:hilbert}
\end{figure}
\begin{figure}[htbp]
 \includevspacePostScript[12.5cm]{RegleConsPeano}{0.95}{0.95}
 \caption{Règles de construction des courbes de Peano}\label{fig:conspeano}
\end{figure}
\begin{figure}[htbp]
 \includevspacePostScript[12.2cm]{Peano}{0.95}{0.95}
 \caption{Courbes de Peano}\label{fig:peano}
\end{figure}
\pagebreak[3]
\begin{corriges}{
\renewcommand{\carrenoir}{\hspace{-4mm}\rule{1.5mm}{1.5mm}\hspace{2.5mm}}
Nous donnons deux programmes, les autres se trouvant dans le dossier {\idt livre/\allowbreak 21\_recursivite}~:
\label{prog:SegmentDeVonKoch}
\input ocaml/SegmentDeVonKoch.tex
\pagebreak[3]
\input ocaml/ConversionDecimalEnBinaire.tex
\renewcommand{\carrenoir}{}
\noindent $\Box$}\end{corriges}
\vfill

%%%%%%%%%%%%%%%%%%%%%%%%%%%%%%%%%%%%%%%%%%%%%%%%%%%%%%%%%%%%%%%%%%%%%%%%%%%
\noindent \framebox[\textwidth]{
\begin{minipage}{\myparboxwidth}
\setlength{\parindent}{2em}
\small \sf
\vspace{\abovedisplayskip}
\section*{Addenda}
 Les langages {\sc algol60} \cite{algol60} puis {\sc lisp} \cite{lisp} furent les premiers à 
 introduire l'usage de la récursivité. La polémique initiale sur l'efficacité des procédures
 récursives est maintenant dépassée, les techniques de compilation étant bien au point. 
 L'idée de récursivité fut à la base de nombreux progrès en informatique, que ce soit les 
 structures de données récursives (comme les arbres), les
 algorithmes récursifs (comme les algorithmes rapides de tri) ou plus récemment, la
 programmation fonctionnelle, dont \ocaml\ est un exemple.

La façon la plus simple de comprendre une procédure 
 récursive est d'imaginer qu'à chaque appel l'ordinateur exécute une nouvelle {\em copie} de la
 procédure. 

 Pour montrer que cette succession d'appels se termine, il suffit de montrer qu'il
 existe une fonction entière de la valeur des paramètres qui est po\-si\-ti\-ve et
 décroit strictement lors d'appels successifs. 
 Par exemple la procédure {\idt segment}
 du programme {\idt SegmentDeVonKoch} (page \pageref{prog:SegmentDeVonKoch}) doit être
 appelée avec un paramètre {\idt n} positif ou nul. Si c'est le cas, les appels récursifs sont de la forme
 {\idt  segment (n - 1)} avec {\idt n} $> 0$ de sorte que lors d'appels successifs les paramètres
 effectifs décroissent strictement en restant positifs. Ceci n'est pas possible indéfiniment.
 Par conséquent la procédure {\idt segment} ne boucle pas quand elle est appelée avec un 
 paramètre initial positif ou nul.

\index{Fonction} Les {\em fonctions} sont aux expressions ce que les procédures sont aux instructions. En
 \ocaml, une déclaration de fonction s'écrit comme une déclaration de procédure, mais son corps est
 une expression dont la valeur est le résultat de la fonction appliquée à ses paramètres. Les procedures
 sont simplement des fonctions dont on s'intéresse pas au résultat, qui est dénoté \texttt{()}, la seule
 valeur de type \texttt{unit}.

\vspace{\abovedisplayskip}
\end{minipage}
}
%%%%%%%%%%%%%%%%%%%%%%%%%%%%%%%%%%%%%%%%%%%%%%%%%%%%%%%%%%%%%%%%%%%%%%%%%%%% 

%%%%%%%%%%%%%%%%%%%%%%%%%%%%%%%%%%%%%%%%%%%%%%%%%%%%%%%%%%%%%%%%%%%%%%%%%%%
\noindent \framebox[\textwidth]{
\begin{minipage}{\myparboxwidth}
\setlength{\parindent}{2em}
\small \sf
\vspace{\abovedisplayskip}
 
 
 Par exemple, une constante est une fonction sans paramètres de sorte 
 qu'une déclaration comme~:
 
\smallskip
\mbox{\blc  \blk\idt  \mcl let\idt \ tailleGrille\ =\ 15}\par 
\smallskip

\noindent peut être remplacée par une déclaration de fonction~:

\smallskip
\mbox{\blc  \blk\idt  \mcl let\idt \ tailleGrille\ ()\ =\ 15}\par 
\smallskip
 
\noindent qui s'utilise en écrivant {\idt tailleGrille ()}, comme {\idt valLgX ()}.
Comme second exemple considérons le calcul de $2^{n}$ avec $n \geq 0$ en observant pour la base
que $2^{0} = 1$ et pour l'induction que $2^{n} = 2 \ast 2^{n-1}$~:

\smallskip
\mbox{\blc  \blk\idt  \mcl let rec\idt \ deuxPuissance\ n\ =}\par 
\mbox{\blc  \blk\blk\idt  \mcl if\idt \ n\ =\ 0\ \mcl then\idt \ 1}\par 
\mbox{\blc  \blk\blk\idt  \mcl else\idt \ 2\ *\ deuxPuissance\ (n\ {-}\ 1)}\par 
\smallskip

%\vspace{\abovedisplayskip}
%\end{minipage}
%}
%%%%%%%%%%%%%%%%%%%%%%%%%%%%%%%%%%%%%%%%%%%%%%%%%%%%%%%%%%%%%%%%%%%%%%%%%%%% 

%%%%%%%%%%%%%%%%%%%%%%%%%%%%%%%%%%%%%%%%%%%%%%%%%%%%%%%%%%%%%%%%%%%%%%%%%%%
%\noindent \framebox[\textwidth]{
%\begin{minipage}{\myparboxwidth}
%\setlength{\parindent}{2em}
%\small \sf
%\vspace{\abovedisplayskip}
\noindent Après cette déclaration, l'expression {\idt deuxPuissance 2} vaut 4 et
l'ex\-pres\-sion {\idt deuxPuissance (deuxPuissance 2)} vaut 16 (le programme
{\idt DeuxPuissance} du dossier {\idt livre/\allowbreak 21\_recursivite} dessine la valeur de
{\idt deuxPuissance 16}). Ici {\mcl if}\ldots{\mcl then}\ldots{\mcl else} choisit entre deux
expressions, et non entre deux instructions.

 En \ocaml\ on peut utiliser les fonctions trigonométriques {\idt sin} (sinus),
 {\idt cos} (cosinus), {\idt tan} (tangente) et {\idt atan} (arc tangente). Les angles sont exprimés en
 radians (un angle de $d$ degrés vaut {\idt d /.\ 180.\ *.\ pi} radians, la constante {\idt pi}
 valant $\pi \simeq 3.14159$). Comme pour les
 fonctions {\idt exp} (exponentielle) et
 {\idt log} (logarithme népérien), le
  paramètre est de type {\idt float} (un entier {\idt n} doit être converti par
  {\idt float\_of\_int n}) et le
résultat est de type {\idt float}. Pour la puissance, on utilise
 {\idt a ** x} qui vaut $a^{x} = e^{x \ln a}$. On pourrait définir soi-même la fonction tangente par~:
 
{
\parindent \parindentprog
\medskip
\mbox{\blc  \blk\idt  \mcl let\idt \ tg\ x\ =\ sin\ x\ /.\ cos\ x}\par 
}

\noindent Quand {\idt cos x} vaut {\idt 0.}, la division ne provoque pas d'erreur~: son résultat est
{\idt infinity} (l'infini) ou {\idt neg\_infinity} (moins l'infini).

\vspace{\abovedisplayskip}
\end{minipage}
}
%%%%%%%%%%%%%%%%%%%%%%%%%%%%%%%%%%%%%%%%%%%%%%%%%%%%%%%%%%%%%%%%%%%%%%%%%%%% 

